% page 229 of the facsimile (image p-228 of the scan)
% block 165.1 x 246.3 mm ; left margin 36.9 mm
\pgc{15.240mm}{0.000mm}{
\l{}
\l{}
\l{}
\l{\cel{53}221}
\l{}
\l{\sou{hodometro}. I. Instrumento horlojo-forma, di qua la indexo quik}
\l{\cel{3}kande lu movesas, indikas la quanto de pazi ed anke la disto}
\l{\cel{3}parkurita da ped-iranto dum ta o ca quanto de tempo. - II.}
\l{\cel{3}instrumento uzata por kontar la quanto de rotaci di roto,}
\l{\cel{3}di mashin-arboro, en ta o ca tempo. - DEFIS.}
\l{}
\l{\sou{hoko}. I. Instrumento kun la extremajo kurva, per qua onu povas}
\l{\cel{3}tirar ulo ad su. - II. Mikra peco kurvoza an qua onu povas}
\l{\cel{3}suspendar ulo. - III. (\sou{metaf}.) Ulo kurvoza quale hoko. - DE.}
\l{}
\l{\sou{hola}\cel{1}! Interjeciono por advokar, por incitar ad halto.}
\l{}
\l{\sou{holdo}. (\sou{navig}.)La parto maxim infra di la internajo di navo. - E.}
\l{}
\l{\sou{holoblastika}. (\sou{biol.}) Dicesas pri la ovo qua su dividas tote}
\l{\cel{3}ed en parti inter-egala (ovo +alecita, sen nutro-provizuri.)}
\l{}
\l{\sou{holografa}.\cel{2}(\sou{pri testamento}). Tote skribita per la manuo di la}
\l{\cel{3}testamentanto. - DEFIS.}
\l{}
\l{\sou{holokausto}. Che la Hebrei antiqua : sakrifiko en qua la viktimo}
\l{\cel{3}esis tote kombustata. - EFIS.}
\l{}
\l{\sou{holometro.}\cel{2}(\sou{matem}.)Instrumento per qua onu mezuras la alteso-}
\l{\cel{3}angulo di punto super horizonto.}
\l{}
\l{\sou{holoturio}. (\sou{zool}.) Genero de +radiarii ekinoderma. - L.\cel{1}\sou{Holo}}
\l{\cel{3}\sou{thuria}. - DEFIRS.}
\l{}
\l{\sou{homo.}.\cel{1}\sou{(zool.}) Mamifero bimanua, qua stacas vertikale, dotita ye}
\l{\cel{3}la fakultato parolar e ye raciono. - EFIS. L.\cel{1}\sou{homo}.}
\l{}
\l{\sou{homajo}. I. (\sou{epoki feudala}). Ago da la vasalo qua su deklaras}
\l{\cel{3}devota a lua sinioro e promisas ad ica servor lu fidele e}
\l{\cel{3}sen restrikto. - II. Ago submisa e respekta. - III. Respekto-}
\l{\cel{3}signo humilesala. - EFI.}
\l{}
\l{\sou{homardo}. (\sou{zool.})\cel{2}Genero de krustacei dek-peda, di qua la du}
\l{\cel{3}gambi avana esas grosega e similesas pinci. L.\cel{1}\sou{homarus}}
\l{\cel{2}vulgaris. - DFI.}
\l{}
\l{\sou{homeopatio}. (\sou{medic.}) Sistemo terapiala - da S. Hahnemann,}
\l{\cel{3}1755-1843) qua konsistas ye traktar la morbi per dozi medi-}
\l{\cel{3}kamentala infinitezima, qui produktas, en la individuo sana,}
\l{\cel{3}la simptomi di morbo analoga ad olta quan onu probas kombatar.}
\l{\cel{3}- DEFIRS.}
\l{}
\l{homeosexuala. (\sou{patol.}) Individuo qua sentas afineso sexuala}
\l{\cel{3}nur kun la personi di sua propra sexuo. - DEF.}
\l{}
\l{homerala. (Rido) nerepresebla (quan Homeros, en la fino di la}
\l{\cel{3}kanto unesma di\cel{1}\sou{Iliado}, atribuas a la dei).}
}
